% Синтетический кейс: BibTeX с двумя библиографиями (multibbl), списки из .bbl,
% пути с ./, второй титульный блок после \maketitle, figure*, <<кавычки>>.
\documentclass[12pt]{article}
\usepackage[T2A]{fontenc}
\usepackage[utf8]{inputenc}
\usepackage[english,russian]{babel}
\usepackage{amsmath}
\usepackage{multibbl}
\newbibliography{rus}
\newbibliography{eng}

\begin{document}
\title{Статья с BibTeX}
\author{А. Автор}
\maketitle

\section{Цитаты}
В АО <<НПК <<Криптонит>> читают~\cite{rus}{knuth2000} и~\cite{rus}{knuth2000,lamport1978}.

\begin{figure*}[h]
\centering
Текст рисунка на две колонки
\caption{Широкий рисунок}\label{fig:wide}
\end{figure*}
См. рис.~\ref{fig:wide}.

\bibliographystyle{rus}{unsrt}
\bibliography{rus}{./bib/refs}{Библиография}

\title{Article with BibTeX}
\author{A. Author}
\nocite{eng}{*}
\bibliographystyle{eng}{unsrt}
\bibliography{eng}{./bib/refs_en}{References}
\end{document}
